% !TEX root = ../main.tex
\chapter{KESIMPULAN DAN SARAN}
\label{chap:kesimpulan}

% ============================================================
\section{Kesimpulan}
\label{sec:kesimpulan}
% ============================================================

Penelitian ini merancang, mengimplementasikan, dan mengevaluasi sistem pelacakan pelanggaran penggunaan helm di Institut Teknologi Sumatera (Itera). Sistem dibangun sebagai purwarupa pada satu komputer menggunakan arsitektur Kappa (\textit{Kappa architecture}) berbasis kontainer Docker, dengan inferensi video dan analitik data besar (\textit{big data analytics}) untuk pengelolaan hasil deteksi. Berdasarkan hasil pengujian, terdapat tiga kesimpulan utama yang menjawab tujuan penelitian.

\begin{enumerate}[1.]
  \item Model deteksi objek berbasis YOLOv8 berhasil dilatih menggunakan dataset kustom CCTV kampus. Model memperoleh \textit{Mean Average Precision} (mAP@50) sebesar 91,4\% pada data validasi dan 90,0\% pada data uji. Nilai \textit{precision} sebesar 87,5\% dan \textit{recall} sebesar 86,5\% hanya berselisih 1,0\%, sehingga kinerja model relatif seimbang antara ketepatan prediksi dan kemampuan menemukan objek. Hasil tersebut menunjukkan bahwa model mampu membedakan kelas \texttt{MOTORCYCLE}, \texttt{DRIVER\_HELMET}, dan \texttt{DRIVER\_NO\_HELMET} pada variasi pencahayaan dan kondisi visual luar ruang.

  \item Integrasi YOLOv8 dan ByteTrack mengubah keluaran deteksi menjadi aliran peristiwa (\textit{event stream}) yang lebih terstruktur. Filter asosiasi spasial (\textit{Intersection over Area}) dan konfirmasi temporal \textit{N-Frame} digunakan untuk menyaring kandidat pelanggaran sebelum data dikirim ke sistem hilir. Pada pengujian 3.001 \textit{frame}, rangkaian filter menyaring 9.544 deteksi mentah menjadi 2 \textit{event} pelanggaran unik yang disimpan ke \textit{data mart}. Pada pemrosesan lokal di satu perangkat, sistem mencatat \textit{throughput} rata-rata 18,9 FPS dan latensi total 53 milidetik pada perangkat pengujian tanpa GPU diskrit. Pengiriman data melalui Apache Kafka mencatat median latensi transmisi 2,65 milidetik, sedangkan Apache Spark Structured Streaming memproses data secara asinkron tanpa membentuk antrean \textit{backpressure}.

  \item \textit{Data mart} berbasis skema bintang (\textit{star schema}) pada PostgreSQL berhasil digunakan sebagai penyimpanan analitik pelanggaran helm. Pada rekaman operasional 10 jam, yaitu pukul 07.00--17.00 WIB, sistem mencatat 113 kejadian pelanggaran. \textit{Dashboard} Grafana menampilkan pola kepadatan bimodal, dengan puncak tertinggi pada pukul 14.45--15.15 WIB yang berkaitan dengan kondisi pascahujan dan transisi jadwal kelas. Dengan rancangan tersebut, sistem tidak hanya berfungsi sebagai alat deteksi, tetapi juga sebagai komponen Sistem Pendukung Keputusan (\textit{Decision Support System}) untuk membantu UPT K3L Itera membaca pola waktu dan lokasi pelanggaran secara berbasis data.
\end{enumerate}

% ============================================================
\section{Saran}
\label{sec:saran}
% ============================================================

Sistem telah berjalan sesuai batasan penelitian, tetapi masih dapat dikembangkan sebelum diterapkan pada skala produksi (\textit{production-grade}). Rekomendasi penelitian lanjutan disusun ke dalam tiga arah pengembangan berikut.

\begin{enumerate}[1.]
  \item {Penerapan pengujian multikamera.}
    Penelitian ini diuji pada satu kamera di gerbang utama. Pengujian berikutnya dapat dilakukan pada beberapa titik penting, seperti gerbang sekunder, persimpangan internal, dan jalur sepeda motor. Penggunaan beberapa kamera perlu disertai mekanisme \textit{re-identification} antarkamera agar \texttt{Track ID} sesi dapat dikaitkan dengan identitas objek lintas kamera yang lebih konsisten. Pengembangan ini penting karena Bab~\ref{chap:hasil} menunjukkan bahwa \texttt{Track ID} ByteTrack masih bersifat identitas sesi, bukan identitas permanen.

  \item {Penerapan MLOps untuk siklus hidup model.}
    Pelatihan model pada penelitian ini dilakukan satu kali menggunakan Roboflow Train 3.0. Pada penggunaan jangka panjang, model perlu diperbarui karena sudut kamera, kondisi cuaca, musim, dan variasi kendaraan dapat berubah. Penelitian lanjutan dapat membangun jalur MLOps yang mencakup registri model (\textit{model registry}), pelacakan versi dataset, dan pelatihan ulang otomatis menggunakan perangkat seperti MLflow, DVC, dan Apache Airflow. Pemantauan pergeseran distribusi (\textit{model drift}) juga dapat dilakukan dengan Evidently AI atau WhyLabs, sehingga pelatihan ulang dapat dipicu ketika metrik turun di bawah ambang tertentu. Selain itu, mekanisme \textit{human-in-the-loop} dapat digunakan untuk memeriksa \textit{event} berkeyakinan rendah dan memperbaiki mutu anotasi.

  \item {Perluasan deteksi ke jenis pelanggaran lalu lintas lain.}
    Sistem saat ini berfokus pada pelanggaran helm. Pada penelitian berikutnya, \textit{pipeline} yang sama dapat dikembangkan untuk mendeteksi pelanggaran lain, seperti melawan arus, motor masuk jalur pejalan kaki, parkir liar di luar petak, atau penerobosan rambu larangan di area kampus. Pelanggaran berbasis pergerakan membutuhkan analisis lintasan (\textit{trajectory}) menggunakan \texttt{Track ID} ByteTrack yang dibandingkan dengan poligon zona dan arah lajur yang ditentukan operator. Pengembangan ini juga memerlukan dataset tambahan karena setiap jenis pelanggaran membutuhkan anotasi dan aturan validasi yang berbeda. Dengan pengembangan tersebut, satu \textit{pipeline} pemrosesan dapat memantau beberapa kategori pelanggaran, sedangkan \textit{dashboard} Grafana cukup diperluas dengan panel tambahan per kategori.
\end{enumerate}

Secara keseluruhan, penelitian ini menghasilkan rancangan sistem analitik video yang menggabungkan YOLOv8, ByteTrack, Kafka--Spark, \textit{data mart}, dan \textit{dashboard} Grafana dalam satu alur kerja. Sistem ini mengubah rekaman CCTV menjadi \textit{event} pelanggaran yang terukur dan dapat dianalisis. Tiga arah pengembangan di atas dapat menjadi dasar untuk memperluas sistem dari purwarupa satu kamera menjadi sistem pemantauan kampus yang lebih siap digunakan secara operasional.
